\cotefile{54}
\scanpage{2}
Dualité en Cohomologie \unsure{des|de} V. A.
% archive sheet number: 1

\scanpage{3}
% Lines 1–2. The two left entries are one key.
\[
\begin{aligned}
 A&\quad\text{V.A.}\\
 B&\quad\text{duale}
\end{aligned}
\qquad
 H^*(A)=\underline{\underline{H}}^*(A,\underline{\underline{\Omega}}_A^*)
 \quad\text{\unsure{Coh.} de Hodge-DeRham}
\]
% Lines 3–4
\[
\begin{aligned}
 t_B&=H^1(A,\underline{O}_A)\\
 t_A&=H^0(A,\Omega_A^1)^{\deleted{\unsure{*}}\vee}\simeq H^n(A,\Omega_A^{n-1})
\end{aligned}
\]

% Lines 5–6
\margin{$\deleted{\unsure{\xi|S}}\unsure{S_A}$}
\diagram{\begin{tikzcd}[ampersand replacement=\&,column sep=large]
 0 \arrow[r,phantom,"{\unsure{\relbar|\to}}" description] \& t_A^\vee \arrow[r] \&
 \deleted{\unsure{H}}\unsure{\mathbb{H}|H}^1(A) \arrow[r] \& t_B \arrow[r] \& 0
\end{tikzcd}}
extension canonique (du moins en car. 0 !)

% Lines 7–9
Autre interprétation : Considérons l'extension \uline{universelle} de $A$ par un groupe vectoriel \inserted{$V$}, \deleted{elle} \inserted{une telle extension \deleted{\unsure{se|est}}} se \deleted{\unsure{fait}} \inserted{classifie} par $H^1(A,\deleted{\unsure{\underline{O}_A|\mathcal{O}_A}}V)=\deleted{V}\inserted{\operatorname{Hom}(H^1(A,O_A)^\vee\unsure{\relbar|,|\to}V)}$, donc l'extension universelle est
% Line 10
\diagram{\begin{tikzcd}[ampersand replacement=\&,column sep=large]
 0 \arrow[r,dash] \& \unsure{\mathbb{V}|W}(\deleted{\unsure{H}}t_B^\vee) \arrow[r,-{Stealth[length=2pt,width=2pt]}] \& E \arrow[r,dash] \& A \arrow[r,-{Stealth[length=2pt,width=2pt]}] \& 0
\end{tikzcd}}
et passant aux espaces tangents on trouve une extension
% Line 13
\diagram{\begin{tikzcd}[ampersand replacement=\&,column sep=large]
 0 \arrow[r] \& t_B^\vee \arrow[r] \& t_E \arrow[r] \& t_A \arrow[r] \& 0
\end{tikzcd}}

% Lines 14–16
Je dis que cette extension n'est autre que $S_B$, (ce qui donne en passant
\[
 t_E=H^1(B) \tag*{)}
\]
% Lines 17–19
et que de plus elle est \uline{transposée} de $S_A$, ce qui donne en passant
\[
 \boxed{H^1(A)\ \text{et}\ H^1(B)\ \text{sont duels l'un de l'autre}}
\]
% Line 20
\margin{$\Vert$}Plus \unsure{gén.,|gén,|fin,} il faut \unsure{présumer|prouver} que $H^*(A)$ et $H^*(B)$
% archive sheet number: 2

\scanpage{4}
% A single diagonal cancellation crosses the typescript and reaches into the handwritten draft below it.
\deleted{\uline{Corollaire} \deleted{\unsure{4.2.|4.3.} quater} Soient $A$ un anneau à spectre irréductible, $\underline{p}$ son unique idéal premier minimal, $B$ une $A$-algèbre de présentation finie $M$ un $B$-module de présentation finie. Supposons que $M_{\underline{p}}$ soit plat\inserted{ sur $A_{\underline{p}}$}, alors il existe $f\handwritten{\in}A-\underline{p}$ tel que $\unsure{\mathbb{M}|M}_f$ soit \uline{libre} sur $A_f$.

\deleted{Supposons d'abord $A$ noethérien, donc $A_{\underline{p}}$ local \uline{artinien}, alors l'hypothèse que $M_{\underline{p}}$ est plat sur $A_{\underline{p}}$} En vertu de \deleted{\unsure{4.1. ter|4.1. bis}}\inserted{Prop.A}, quitte à remplacer $A$ par un $A_f$ \deleted{\unsure{sur}} avec $f\quad A-\underline{p}$ convenable, on peut supposer \unsure{qu'i|qu'il} existe un idéal $J$ \inserted{de type fini} dans $A$ tel que, pour toute $A$-algèbre $A'$, on ait $M\otimes_A A'$ est libre sur $A'$, ou encore plat sur $A'$, sss $A\quad A'$ s'annule sur $J$. Appliquant ceci au cas $A'=A_{\underline{p}}$, on voit que l'on a $JA_{\underline{p}}=0$ i.e. $J_{\underline{p}}=0$. Comme $J$ est de type fini, il s'ensuit qu'il existe $f\quad A-\underline{p}$ tel que $J_f=0$ i.e. $JA_f=0$, ce qui implique que $M_f$ est libre sur $A_f$, cqfd

Je m'aperçois à l'instant d'une amélioration dans une autre direction de 4.1. , dont je vais reprendre l'énoncé:

\uline{Proposition} 4.1. Soient $A$ un anneau, $I$ un idéal \deleted{nilpotent} de $A$, $(A_i)_{i\;I}$ une famille de $A$-algèbres telle que l'intersection des noyaux des homomorphismes $A\quad A_i$ soit réduite à 0, $M$ un $A$-module. O\deleted{n}n suppose $M\otimes_A A/I$ libre sur $A/I$ (resp. plat sur $A/I$, \uline{et} l'ensemble d'indices $i$ fini), et pour tout $i$, $M\otimes_A A_i$ libre (resp. plat) sur $A_i$. Alors \inserted{\handwritten{$\forall n$,}} \inserted{\handwritten{$M\otimes_A(A/I^{n+1})$}} $M$ est \deleted{plat sur \unsure{$A$.|$A$}} libre (resp. plat) sur $A\inserted{\handwritten{/I^{n+1}}}$.

Dans le cas non respé, il suffit de prouver aussi la platitude de $M_{\inserted{\handwritten{n}}}$ sur $A_{\inserted{\handwritten{n}}}$, qui entraine la liberté, grâce à \inserted{\handwritten{$M_0=$}}$M\otimes_A A/I$ libre et $I$ nilpotent.

% Successive starts at a handwritten corollary, crossed by a second, looping stroke.
\handwritten{\uline{Corollaire} \deleted{Supposons $A$ anneau \uline{noethérien}\inserted{\unsure{noeth.}}, $M$ un $A$-module,} \margin{$\Bigg[$}\deleted{$(\unsure{\mathfrak{a}|\mathfrak{U}|I}_i)_{i\in\unsure{I}}$ des idéaux de $A$, \inserted{$(\bigcap\mathfrak{a}_i=0)$} \unsure{Supposons} que $\forall i$, $M\otimes_A A/\mathfrak{a}_i$ soit \unsure{plt|plat} sur $\mathfrak{a}_i$. Alors pour tout entier $n$,} Supposons \deleted{\unsure{que}} qu'on est dans l'un des cas \unsure{suivants}

(i) $I$ nilpotent}}%

\handwritten{(ii) $A$ \unsure{noeth.} et $M$ \unsure{idél\textsuperscript{emt}|idéalement|idéliquement} séparé pour $I$.

Alors \deleted{\unsure{les deux conditions sont}}}
\[
\handwritten{M\ \unsure{\text{plt}|\text{plat}}/A
 \iff
 \begin{aligned}
 M_0&\ \text{plt sur }A_0\\
 M_i&\ \text{plt sur }A_i,\ \forall i
 \end{aligned}}
\]
% archive sheet number: 3

\scanpage{5}
% Lines 1–2, continuation of the final line of scan 3.
\margin{$\Vert$}sont duels l'un de l'autre, et que $H^p(A,\Omega_A^q)$ est dual de $H^q(B,\Omega_A^p)$.

% Lines 3–9. There is an opening square bracket in the left margin.
[ Ce dernier fait, grâce à $H^*(A,\Omega_A^*)=\bigwedge\bigl(H^*(A,\Omega_A^{*})\bigr)^{(1)}$ se ramène au cas de la \unsure{dim 1|dim 1.}, i.e. par symétrie à la dualité de $H^{\deleted{\unsure{0}}1}(A,\deleted{\unsure{\Omega}}\underline{O}_A)$ et $H^0(B,\deleted{\unsure{\underline{O}}}\Omega_B^1)$, qui est immédiate. \deleted{De plus, \unsure{on},} pour avoir une dualité entre $H^*(A)$ et $H^*(B)$, il \unsure{suffit|\deleted{suffit}|suffirait|\deleted{suffirait}} grâce à la structure connue de la définir entre $H^1(A)$ et $H^1(B)$, ce qui ne saurait offrir de difficultés....

% Lines 10–15
Il y a lieu enfin de comparer ces extensions et dualités avec ce que donne la \uline{dualité de Cartier}. \unsure{Envisageons|Considérons} (\unsure{la car.|le cas} étant $p>0$) $ {}_pA$ qui est dual au sens Cartier de $ {}_pB$. On a aussi $ {}_F A=\operatorname{Gr}(t_A)\subset {}_p A$\unsure{.|-}\,; de degrés respectifs $p^{\deleted{\unsure{2|2d}}d}$ et $p^{2d}$, où $d=\dim A$.
% Lines 16–17
\[
 \left.\begin{aligned}
 {}_F A&\subset {}_p A\\
 {}_F B&\subset {}_p B
 \end{aligned}\right]\ \text{duels l'un de l'autre}
\]
% Lines 18–19
Je dis que $ {}_F A$ \uline{et} $ {}_F B$ \uline{sont orthogonaux l'un de l'autre}. \deleted{Noter que $ {}_F B$} On le
% archive sheet number: 4

\scanpage{6}
% Typed lines 1–6
\deleted{\uline{Corollaire} 4.1. ter} \inserted{\handwritten{Proposition $A$}} Soit $f:X\quad Y$ un morphisme de présentation finie, avec $Y$ \deleted{\unsure{réduit|affine}}\inserted{\handwritten{$\ne\varnothing$}}, $F$ un Module de présentation finie sur $X$. Alors, quitte à remplacer $Y$ par un ouvert non vide convenable, on a ce qui suit: il existe un sous-schéma fermé \deleted{\unsure{Z d|et d}} $Z$ de $Y$, \inserted{\handwritten{de présentation finie sur $Y$,}} tel que pour tout changement de base $Y'\handwritten{\to}Y$, la condition nécessaire et suffisante pour que $F'$ soit plat par rapport à $Y'$ est que $Y'\quad Y$ se factorise \deleted{par} en $Y'\quad Z\quad Y$.
% Typed lines 7–10
De plus,\inserted{si $Y$ est affine et $X$ est réunion finie d'ouverts affines $U_i$, \unsure{alor|alors}}(quitte à remplacer $Y$ par un ouvert non vide plus petit), on a ceci: si la condition précédente est vérifiée pour $Y'$, et si $Y'$ est affine\inserted{ d'anneau $A'$}, alors \deleted{pour tout ouvert affine $U$ de $\unsure{X'|X}$,}\inserted{pour tout $i$,} $\handwritten{\Gamma}(\unsure{\mathbb{U}|U}_i,F')$ est un $A'$-module \uline{libre}.

% Typed lines 11–17
On peut évidemment supposer $Y$ affine, i.e. $Y=\operatorname{Spec}(A)$. Notant que la situation $(X,F)$ sur $Y$ provient par changement de base d'une situation analogue $(X_{\unsure{o|0}},F_o)$ sur un $Y_o=\operatorname{Spec}(A_o)$, où $A_o$ est un sous-anneau noethérien de $A$, on est ramené trivialement au cas où $A$ est noethérien. Soit $I$ l'idéal des éléments nilpotents de $A$, $A_o=A/I$, $Y_o=\operatorname{Spec}(A_o)$, et dénotons par un $ {}_o$ l'effet du changement de base $Y_o\quad Y$. En vertu du théorème de platitude générique (par 9), quitte à remplacer $Y$ par un ouvert non vide \deleted{\unsure{v|o}}convenable, on peut supposer que les $M_{io}=I(U_{i\;o},F_{i\;o})=M_i\otimes_A A_o$ (où $M_i=I(U_i,F)$ ) sont des $A_o$-modules \uline{libres}. A partir de là, l'assertion résulte facilement de 4.1. bis : pour tout $i$, on trouve un idéal $J_i$ dans $A$, tel que pour \unsure{tout|toute} $A$-algèbre $A'$, $M_i\otimes_A A'$ soit libre, ou encore plat, sur $A'$, sss $A\quad A'$ est nul sur $J_i$. Soit $J$ l'idéal somme des $J_i$, alors pour tout $i$, $M_i\otimes_A A'$ est libre sur $A'$, ou encore plat sur $A'$, sss $A\quad A'$ s'annule sur $J$, i.e. $A\quad A'$ se factorise en $A\quad A/J\quad A'$. Il suffit maintenant de pr\,e\deleted{\unsure{d|n}}ndre $Z=V(J)=\operatorname{Spec}(A/J)$, notant que \deleted{\unsure{la co|le mo}} pour vérifier que $Z$ satisfait à la condition de l'énoncé, il suffit de le vérifier pour des $\deleted{\unsure{X|Z}}\unsure{Y|Z}'$ affines. \deleted{Noter d'ailleurs}
% archive sheet number: 5

\scanpage{7}
% Lines 1–2. The faint reversed typescript shows through from the other face.
formellement sur le \inserted{\unsure{un|$\sim$}} diagramme de deux suites transposées l'une de l'autre
% Lines 3–4
\diagram{\begin{tikzcd}[ampersand replacement=\&,column sep=large]
 \unsure{A|A_1} \arrow[r,"F"] \& A^{(p)}
 \arrow[r,phantom,"{\overset{\unsure{\mathbb{V}|V}}{\unsure{\relbar|\to}}}" description] \& V\\
 B \& B^{(p)} \arrow[l,"V"'] \& B \arrow[l,"F"']
\end{tikzcd}}

% Lines 5–7
On \unsure{trouve|a donc} ainsi la suite exacte
\diagram{\begin{tikzcd}[ampersand replacement=\&,column sep=large]
 \unsure{0|\cdot} \arrow[r,dash] \& \operatorname{Gr}(\underline{t}_A) \arrow[r,-{Stealth[length=2pt,width=2pt]}] \& {}_p A \arrow[r,-{Stealth[length=2pt,width=2pt]}] \& D(\operatorname{Gr}(\underline{t}_B)) \arrow[r] \& 0
\end{tikzcd}}
extension canonique de $D({}_F\deleted{\unsure{A}}B)$ par $ {}_F A$.

% Lines 8–12. The enclosed interline has a point descending after “abéliennes”.
Noter d'ailleurs que $S_A$ est une extension de $p$-alg. de Lie abéliennes \inserted{$t_A^\vee$ \unsure{étant|est} muni de la puiss. \unsure{$p$.ième|$p$-ième} nulle.}, et correspond donc en fait à une extension de groupes \unsure{radiciels}. Cette extension peut-elle se déduire de la connaissance de $ {}_p A$ ??
% archive sheet number: 6

\scanpage{8}
% Typescript upside down in the archive scan.
\mspage{- 19 -}
n\textsuperscript{o} 391

terminer la démonstration, soient $y$ un point de $C$, $z$ un point de $C'=\bigcap\limits_{H\in\underline{\unsure{M|H}}}D_H(C)$. Comme $\underline{H}$ est localement fini, le segment fermé d'extrémités $y$, $z$ ne peut rencontrer qu'un nombre fini d'hyperplans de $\underline{H}$ ; il n'y a par suite qu'un nombre fini d'hyperplans $H_1,..,H_q$ de $\underline{H}$ tels que $z\notin D_{H_i}(C)$ pour $1\leq i\leq q$. Par définition de $C'$, les $H_i$ ne sont pas des murs de $C$, donc, par récurrence sur $i$, on voit que $C$ est l'intersection des $D_H(C)$ pour les $H\in\underline{H}$ \uline{distincts des} $H_i$. Cela prouve que $z\in C$, donc $C'\subset C$, et par suite $C'=C$.

COROLLAIRE 2.- \uline{Soit} $\underline{H}$ \uline{un ensemble localement fini d'hyperplans affines}. \uline{Alors tout hyperplan} $H\in\underline{H}$ \uline{est mur mitoyen de deux chambres déterminées par} $\underline{H}$.

Il suffit d'appliquer la prop. 10 à l'ensemble $\underline{H}-\{H\}$.

COROLLAIRE 3.- \uline{Soient} $\underline{H}$ \uline{un ensemble localement fini d'hyperplans affines,} $C$ \uline{une chambre déterminée par} $\underline{H}$, $H_{\unsure{o|0}}$ \uline{un mur de} $C$. \uline{L'intérieur de} $\bar{C}\cap H_o$ \uline{dans} $H_o$ \uline{est une chambre dans} $H_o$, \uline{formée des points de} $\bar{C}\cap H_o$ \uline{qui n'appartiennent à aucun hyperplan} $H\ne H_o$ \uline{de} $\underline{H}$ ; \uline{ce sont aussi les points} $x$ \uline{de} $\bar{C}\cap H_o$ \uline{qui sont adhérents à} $C$ \uline{et à une seule des autres chambres déterminées par} $\underline{H}$, \uline{et cette chambre} $C_1$ \uline{est mitoyenne à} $C$ \uline{et est la même pour tous les points} $x$ \uline{intérieurs à} $\bar{C}\cap H_o$.

Le fait que $H_o$ soit un mur de $C$ entraîne que $C_o=\bigcap\limits_{H\in\underline{H},H\ne H_o}D_H(C)$ est distincte de $C$ (cor. 1) ; $C$ est donc l'une des deux chambres notées $C'$, $C''$ dans l'énoncé de la prop. 10 ; supposons par exemple que $C=C'$ ; avec les mêmes notations, l'intérieur de $\bar{C}\cap H_o$ dans $H_o$ est $\Gamma$, et on a vu dans la prop. 10 que les points de $\Gamma$ ne sont adhérents qu'à deux des chambres déterminées par $\underline{H}$, $C'$ et la chambre mitoyenne $C''$. Au contraire, si un point $x\in\bar{C}\cap H_o=\bar{\Gamma}$ n'appartient pas à $\Gamma$, il appartient à un des hyperplans $H\ne H_o$ de $\underline{H}$ ; il est
% archive sheet number: 7
